\newcommand{\titulus}{\nomenFesti{S. Wenceslai, Ducis \& Martyris.}
\dies{Die 28. Septembris.}}
\newcommand{\oratio}{\pars{Oratio.}

\noindent Deus, qui beátum mártyrem Vencesláum cælésti regno terrénum postpónere docuísti, eius précibus concéde, ut, nosmetípsos abnegántes, tibi toto corde adhærére valeámus.

\pars{Pro pace in universo mundo.} \scriptura{Sir. 50, 25; 2 Esdr. 4, 20; \textbf{H416}}

\vspace{-4mm}

\antiphona{II D}{temporalia/ant-dapacemdomine.gtex}

\vfill

\noindent Deus, a quo sancta desidéria, recta consília et iusta sunt ópera: da servis tuis illam, quam mundus dare non potest, pacem; ut et corda nostra mandátis tuis dédita, et hóstium subláta formídine, témpora sint tua protectióne tranquílla.

\noindent Per Dóminum nostrum Iesum Christum, Fílium tuum, qui tecum vivit et regnat in unitáte Spíritus Sancti, Deus, per ómnia sǽcula sæculórum.

\noindent \Rbardot{} Amen.}
\newcommand{\invitatorium}{\pars{Invitatorium.}

\vspace{-4mm}

\antiphona{E}{temporalia/inv-regemmartyrumsimplex.gtex}}
\newcommand{\hymnusmatutinum}{\pars{Hymnus.}

\antiphona{VIII}{temporalia/hym-DeusTuorum.gtex}}
\newcommand{\matversus}{\noindent \Vbardot{} Díriget Dóminus mansuétos in iudício.

\noindent \Rbardot{} Docébit mites vias suas.}
\newcommand{\lectioi}{\pars{Lectio I.} \scriptura{Phil. 1, 12-26}

\noindent De Epístola beáti Pauli apóstoli ad Philippénses.

\noindent Scire vos volo, fratres, quia, quæ circa me sunt, magis ad proféctum venérunt evangélii, ita ut víncula mea manifésta fíerent in Christo in omni prætório et in céteris ómnibus, et plures e frátribus in Dómino confidéntes vínculis meis abundántius audére sine timóre verbum loqui. Quidam quidem et propter invídiam et contentiónem, quidam autem et propter bonam voluntátem Christum prǽdicant; hi quidem ex caritáte sciéntes quóniam in defensiónem evangélii pósitus sum, illi autem ex contentióne Christum annúntiant, non sincére, existimántes pressúram se suscitáre vínculis meis. Quid enim? Dum omni modo, sive sub obténtu sive in veritáte, Christus annuntiétur, et in hoc gáudeo; sed et gaudébo, scio enim quia hoc mihi provéniet in salútem per vestram oratiónem et subministratiónem Spíritus Iesu Christi, secúndum exspectatiónem et spem meam quia in nullo confúndar, sed in omni fidúcia sicut semper et nunc magnificábitur Christus in córpore meo, sive per vitam sive per mortem.

\noindent Mihi enim vívere Christus est et mori lucrum. Quod si vívere in carne, hic mihi fructus óperis est, et quid éligam ignóro. Coártor autem ex his duóbus: desidérium habens dissólvi et cum Christo esse, multo magis mélius; permanére autem in carne magis necessárium est propter vos. Et hoc confídens scio quia manébo et permanébo ómnibus vobis ad proféctum vestrum et gáudium fídei, ut gloriátio vestra abúndet in Christo Iesu in me per meum advéntum íterum ad vos.}
\newcommand{\responsoriumi}{\pars{Responsorium 1.} \scriptura{\Rbardot{} 1 Cor. 15, 10 \Vbardot{} 2 Tim. 4, 7; \textbf{H288}}

\vspace{-5mm}

\responsorium{VIII}{temporalia/resp-gratiadeisumidquodsum-CROCHU.gtex}{}}
\newcommand{\lectioii}{\pars{Lectio II.} \scriptura{(Edit. M. Weingart in Svatováclavský sborník, I, Pragæ 1934, pag. 974-980)}

\noindent E Legénda prima palæosláva.

\noindent Cum pater eius Vratisláus mórtuus esset, constituérunt Bohémi Vencesláum príncipem. Et Dei grátia in fide perféctus erat. Omnibus enim paupéribus bene faciébat, nudos vestiébat, esuriéntes alébat, peregrínos excipiébat secúndum vocem evangélicam. Víduis non patiebátur iniúriam fíeri, hómines omnes, páuperes et dívites amábat, Deo serviéntibus ministrábat, ecclésias multas ornábat.

\noindent Supérbi vero facti sunt viri Bohémi et persuasérunt Bolesláo, fratri eius minóri, dicéntes: «Occisúrus te est frater Vencesláus conspírans cum matre et viris suis».}
\newcommand{\responsoriumii}{\pars{Responsorium 2.} \scriptura{\textbf{Codice Raygradensis 17, folio 201r}}

\vspace{-5mm}

\responsorium{I}{temporalia/resp-wencezlausduxgratiae.gtex}

\rubrica{vel ad libitum:}

\vspace{3mm}

\pars{Responsorium 2.} \scriptura{\Rbardot{} Os. 14, 6 \Vbardot{} Ps. 91, 13; \textbf{H373}}

\vspace{-5mm}

\responsorium{I}{temporalia/resp-justusgerminabit-CROCHU.gtex}{}}
\newcommand{\lectioiii}{\pars{Lectio III.}

\noindent Cum essent festivitátes dedicatiónum ecclesiárum in ómnibus civitátibus, Vencesláus visitábat omnes civitátes. Ingréssus est ígitur Bolesláviam civitátem die domínica, in festo Cosmæ et Damiáni. Audíta missa, vóluit Pragam revérti. Bolesláus autem retínuit eum scelésta mente dicens: «Cur abitúrus es, frater?». Mane facto campánam pulsavérunt ad offícium matutínum. Vencesláus vero, audíto campánæ sono, dixit: «Laus tibi, Dómine, qui dedísti mihi vívere ad hoc mane». Et surréxit et ivit ad offícium matutínum.

\noindent Statim assecútus est eum Bolesláus in iánua. Vencesláus respéxit ad eum et dixit: «Frater, bonus eras nobis fámulus heri». Bolesláo autem diábolus inclinávit se ad aures et pervértit cor eius et, evagináto gládio, respóndit ei dicens: «Nunc volo tibi mélior fíeri». His dictis, ferit caput eius gládio.

\noindent Vencesláus vero convérsus ad eum dixit: «Quid in ánimo habes, frater?». Et prehénsum humi eum prostrávit. Accúrrit autem quidam de consiliáriis Boleslái et percússit manum Venceslái. Hic manu vulnerátus, fratre dimísso, confúgit ad ecclésiam. Maléfici autem duo occidérunt eum in ecclésiæ iánua. Alius accúrrens, latus eius transfódit gládio. Vencesláus inde statim efflávit ánimam cum his verbis: \emph{In manus tuas, Dómine, comméndo spíritum meum.}}
\newcommand{\responsoriumiii}{\pars{Responsorium 3.} \scriptura{\textbf{Codice Raygradensis 17, folio 201v}}

\vspace{-5mm}

\responsorium{III}{temporalia/resp-noctesurgensagrum.gtex}{}

\rubrica{vel ad libitum:}

\vspace{3mm}

\pars{Responsorium 3.} \scriptura{\textbf{H373}}

\vspace{-5mm}

\responsorium{III}{temporalia/resp-istecognovit-CROCHU-cumdox.gtex}{}}
\newcommand{\hymnuslaudes}{\pars{Hymnus.}

\cuminitiali{VII}{temporalia/hym-DiesVenit.gtex}}
\newcommand{\laudes}{\pars{Psalmus 1.} \scriptura{\textbf{R17 205r}}

\vspace{-4mm}

\antiphona{I g}{temporalia/ant-laudemotus-cgp.gtex}

%\vspace{-2mm}

\scriptura{Psalmus 62.}

%\vspace{-1mm}

\initiumpsalmi{temporalia/ps62-initium-i-g.gtex}

%\vspace{-1.5mm}

\input{temporalia/ps62-i-g.tex} \Abardot{}

\vfill
\pagebreak

\pars{Psalmus 2.} \scriptura{\textbf{R17 206r}}

\vspace{-4mm}

\antiphona{IV E}{temporalia/ant-morbidisrefugium-cgp.gtex}

\vspace{-2mm}

%\trAntIV

\scriptura{Canticum trium puerorum, Dan. 3, 57-88 et 56}

\vspace{-2mm}

\initiumpsalmi{temporalia/dan3-initium-iv-E.gtex}

\input{temporalia/dan3-iv-E-sinedox.tex}

\rubrica{Hic non dicitur Gloria Patri, neque Amen.}

\vfill

\antiphona{}{temporalia/ant-morbidisrefugium-cgp.gtex}

\vfill
\pagebreak

\pars{Psalmus 3.} \scriptura{\textbf{R17 206r}}

\vspace{-4mm}

\antiphona{V a}{temporalia/ant-laudantdominum-cgp.gtex}

%\vspace{-2mm}

\scriptura{Psalmus 149}

%\vspace{-2mm}

\initiumpsalmi{temporalia/ps149-initium-v-a-auto.gtex}

\input{temporalia/ps149-v-a.tex} \Abardot{}

\vfill
\pagebreak}
\newcommand{\lectiobrevis}{\pars{Lectio Brevis.} \scriptura{Sir. 39, 6.8.12.13-14}

\noindent Cor suum tradet ad vigilándum dilúculo ad Dóminum, qui fecit illum; spíritu intellegéntiae replébitur. Collaudábunt multi sapiéntiam eius; non recédet memória eius, et nomen eius requirétur a generatióne in generatiónem. Sapiéntiam eius enarrábunt gentes, et laudem eius enuntiábit ecclésia.}
\newcommand{\responsoriumbreve}{\pars{Responsorium breve.} \scriptura{Ps. 20, 6}

\cuminitiali{VI}{temporalia/resp-magnaestgloria.gtex}}
\newcommand{\preces}{\noindent \noindent Fratres, Salvatórem nostrum, testem fidélem,~\gredagger{} per mártyres interféctos propter verbum Dei,~\grestar{} celebrémus, clamántes:

\Rbardot{} Redemísti nos Deo in sánguine tuo.

\noindent Per mártyres tuos, qui líbere mortem in testimónium fídei sunt ampléxi,~\grestar{} da nobis, Dómine, veram spíritus libertátem.

\Rbardot{} Redemísti nos Deo in sánguine tuo.

\noindent Per mártyres tuos, qui fidem usque ad sánguinem sunt conféssi,~\grestar{} da nobis, Dómine, puritátem fideíque constántiam.

\Rbardot{} Redemísti nos Deo in sánguine tuo.

\noindent Per mártyres tuos, qui, sustinéntes crucem, tua vestígia sunt secúti,~\grestar{} da nobis, Dómine, ærúmnas vitæ fórtiter sustinére.

\Rbardot{} Redemísti nos Deo in sánguine tuo.

\noindent Per mártyres tuos, qui stolas suas lavérunt in sánguine Agni,~\grestar{} da nobis, Dómine, omnes insídias carnis mundíque devíncere.

\Rbardot{} Redemísti nos Deo in sánguine tuo.}
\newcommand{\benedictus}{\pars{Canticum Zachariæ.}

\vspace{-6mm}

\antiphona{V a}{temporalia/ant-cordeetlingua-cgp.gtex}

\vspace{-3mm}

\scriptura{Lc. 1, 68-79}

\vspace{-2mm}

\initiumpsalmi{temporalia/benedictus-initium-v-a.gtex}

\vspace{-1.5mm}

\input{temporalia/benedictus-v-a.tex} \Abardot{}}
\newcommand{\benedicamuslaudes}{\cuminitiali{II}{temporalia/benedicamus-solemnism-laud.gtex}}
\include{hebdomadaxxvi}
\include{feriaii}
